% =====================================================================
% ADICOES AO desafio_natura.tex (Luiz Salles-Neto, set/2026)
%
% Cada bloco abaixo indica ONDE entra no documento principal.
% Este arquivo compila sozinho (preambulo proprio) para conferencia
% visual, mas o objetivo e' copiar bloco a bloco para o Overleaf.
%
% Notacao: usa k para ciclo e sigma para shape porque s (setor),
% d (slot), a (dia) e c (CD) ja estao ocupados no documento principal.
% =====================================================================
\documentclass[11pt]{article}

\usepackage[utf8]{inputenc}
\usepackage[T1]{fontenc}
\usepackage[brazilian]{babel}
\usepackage{amsmath,amssymb}
\usepackage{booktabs}
\usepackage{geometry}
\usepackage{hyperref}
\geometry{margin=1in}

\title{Adições ao \texttt{desafio\_natura.tex}\\[4pt]
\large Previsão da demanda: definição de $L_{s,k}$, candidatos de modelo
e critério de comparação}
\author{}
\date{}

\begin{document}
\maketitle

\noindent Este documento reúne o que precisa ser acrescentado ao paper
principal para que o parâmetro $q_{s,a,d}$ da Seção~3.2 deixe de ser um
dado sem origem e passe a ser a saída de um módulo de previsão
especificado. Os blocos estão na ordem em que aparecem no documento
principal.

% =====================================================================
% BLOCO A  -->  SECAO 3.1 (Conjuntos): acrescentar duas linhas na tabela
%               e emendar a descricao de A.
% =====================================================================
\section*{Bloco A — §3.1 Conjuntos}

Acrescentar duas linhas à tabela de conjuntos:

\begin{center}
\begin{tabular}{@{}lll@{}}
\toprule
Símbolo & Descrição & Domínio \\
\midrule
$K$ & conjunto de ciclos (campanhas) do horizonte & $k \in \{1,\dots,|K|\}$ \\
$t$ & índice de dia dentro da janela de vendas & $t \in \{0,\dots,W_k-1\}$ \\
\bottomrule
\end{tabular}
\end{center}

\noindent E emendar a descrição de $A$, hoje ``conjunto de dias do
horizonte de simulação (dias-calendário)'', para explicitar \emph{quais
ciclos} esse horizonte cobre:

\begin{quote}
$A$ — conjunto de dias-calendário do horizonte de simulação,
correspondente aos ciclos $k \in K$.
\end{quote}

\noindent Essa emenda não é cosmética: a Seção~4.1.5 estima $|A| \approx
30$ dias, o que corresponde a aproximadamente uma campanha. Se $|K| = 1$,
o horizonte de previsão é de um ciclo à frente, e isso precisa estar
escrito antes da Seção de validação (Bloco~C.6).

% =====================================================================
% BLOCO B  -->  SECAO 3.2 (Parametros): corrigir q, acrescentar 5 linhas,
%               acrescentar a equacao de decomposicao e duas ressalvas.
% =====================================================================
\section*{Bloco B — §3.2 Parâmetros}

\subsection*{B.1 Correção da unidade de $q_{s,a,d}$}

A descrição atual diz ``quantidade de \textbf{pedidos} prevista pelo setor
$s$ no dia $a$''. A base de demanda disponível não tem coluna de pedidos:
a única coluna de quantidade é \texttt{total\_volumes\_mascarado}.
Substituir por:

\begin{quote}
$q_{s,a,d}$ — \textbf{volume} previsto do setor $s$ no dia $a$, dado que
seu ciclo iniciou em $d$.
\end{quote}

\noindent E acrescentar, logo abaixo da tabela de parâmetros:

\begin{quote}
$q_{s,a,d}$ e $\mathrm{Cap}_{c,a}$ são expressos na mesma unidade
(volumes) e na mesma escala da base mascarada. Os valores de volume da
base são mascarados; a restrição (3.5.4) só é válida se a capacidade
estiver na mesma escala.
\end{quote}

\subsection*{B.2 Linhas novas na tabela de parâmetros}

\begin{center}
\begin{tabular}{@{}lll@{}}
\toprule
Símbolo & Descrição & Domínio \\
\midrule
$L_{s,k}$ & volume total previsto do setor $s$ no ciclo $k$ & $\ge 0$ \\
$\sigma_{s,t}$ & fração do volume do setor $s$ no dia $t$ da janela &
  $\ge 0$, $\sum_t \sigma_{s,t} = 1$ \\
$W_k$ & duração, em dias, da janela de vendas do ciclo $k$ & $\ge 1$ \\
$a^{\mathrm{ini}}_k$ & dia-calendário de início do ciclo $k$ & $\in A$ \\
$\delta_d$ & deslocamento, em dias, da abertura da janela do slot $d$ & $\ge 0$ \\
\bottomrule
\end{tabular}
\end{center}

\subsection*{B.3 Equação de decomposição}

Acrescentar como equação numerada (3.2.1), logo após a tabela de
parâmetros:

\begin{equation}
q_{s,a,d} =
\begin{cases}
L_{s,k}\,\sigma_{s,t}, & t = a - a^{\mathrm{ini}}_k - \delta_d,
  \quad 0 \le t < W_k, \\[4pt]
0, & \text{caso contrário.}
\end{cases}
\end{equation}

\noindent Ou seja: o parâmetro $q_{s,a,d}$ é o produto de um \emph{escalar
de volume} por ciclo, $L_{s,k}$, por uma \emph{curva de distribuição
intra-ciclo}, $\sigma_{s,t}$, deslocada no calendário conforme o slot
$d$ atribuído ao setor. As duas frentes de previsão são separadas
exatamente por essa equação: $L_{s,k}$ é o objeto da Seção~4 deste
documento, $\sigma_{s,t}$ é objeto de trabalho paralelo.

\subsection*{B.4 Duas ressalvas a declarar}

Acrescentar, logo após a equação:

\begin{quote}
\textbf{Hipótese (invariância do volume à atribuição).} $L_{s,k}$ não
depende de $d$: a combinação Bloco/Subloco desloca a demanda do setor no
calendário, mas não altera seu volume total no ciclo. Essa hipótese já
está implícita na forma como $q_{s,a,d}$ foi parametrizado, e é o que
permite prever $L_{s,k}$ sem conhecer a atribuição. Ela precisa ser
verificada empiricamente (Seção~4.7): se a janela de um slot for truncada
por fins de semana ou feriados conforme o dia de abertura, o volume total
pode depender de $d$, e nesse caso $\sigma$ passa a ser indexado também
por $d$.

\textbf{Previsão pontual.} $q_{s,a,d}$ é uma previsão pontual, sem
intervalo de incerteza associado. As restrições (3.5.2)--(3.5.4) tratam
$q$ como determinístico.
\end{quote}

% =====================================================================
% BLOCO C  -->  SECAO NOVA 4 "Previsao da demanda".
%               O Simulated Annealing passa a ser a Secao 5.
% =====================================================================
\section*{Bloco C — Nova §4 ``Previsão da demanda''}

O conteúdo abaixo entra como uma seção nova, posicionada \emph{antes} da
Estrutura do Simulated Annealing (que passa a ser a Seção~5), já que o SA
consome $q_{s,a,d}$ como entrada. As referências cruzadas a equações por
número (3.5.2, 3.5.4 etc.) continuam válidas após a renumeração.

\subsection*{4.1 Base de dados}

A base histórica de demanda tem uma linha por setor e período, com as
colunas:

\begin{center}
\begin{tabular}{@{}ll@{}}
\toprule
Coluna & Conteúdo \\
\midrule
\texttt{cd\_setor} & identificador do setor ($s$) \\
\texttt{aa\_ciclo} & ano do ciclo \\
\texttt{nm\_ciclo} & número do ciclo dentro do ano \\
\texttt{mes\_pedido} & mês em que o pedido foi registrado \\
\texttt{semana\_pedido} & semana em que o pedido foi registrado \\
\texttt{total\_volumes\_mascarado} & volume, em escala mascarada \\
\bottomrule
\end{tabular}
\end{center}

\noindent O par $(\texttt{aa\_ciclo}, \texttt{nm\_ciclo})$ identifica o
ciclo $k$; o período sazonal a ser usado pelos modelos sazonais da
Seção~4.4 é $\max(\texttt{nm\_ciclo})$, o número de ciclos por ano.

\paragraph{Extração de $L_{s,k}$.} O alvo de previsão é obtido por
agregação direta:
\[
L_{s,k} \;=\; \sum_{\text{linhas do setor } s \text{ no ciclo } k}
\texttt{total\_volumes\_mascarado}.
\]

\paragraph{O que a base não contém.} Nenhuma das informações a seguir
está nessa base, e todas são necessárias para fechar a formulação:
a partição $S_c$ (qual CD atende cada setor); a atribuição As-Is
$x^{\mathrm{AsIs}}_{s,d}$ (em qual Bloco/Subloco cada setor esteve em cada
ciclo); as datas de início de ciclo $a^{\mathrm{ini}}_k$; e as capacidades
$\mathrm{Cap}_{c,a}$. A ausência de $x^{\mathrm{AsIs}}_{s,d}$ histórico é a
mais restritiva: sem ela não é possível estimar $\sigma_{s,t}$, que é
definida em relação ao dia de abertura da janela, nem testar a hipótese de
invariância do Bloco~B.4.

\subsection*{4.2 Definição e nível de agregação}

$L_{s,k}$ é um \textbf{escalar por par (setor, ciclo)} --- um valor por
setor por campanha, não uma série diária. A distribuição do volume ao
longo dos dias do ciclo é inteiramente responsabilidade de
$\sigma_{s,t}$. Formalmente, $L_{s,k}$ é o total efetivamente observado
do setor durante a janela de vendas do ciclo, o que corresponde à
agregação da Seção~4.1 e é consistente com a equação~(3.2.1):
\[
L_{s,k} \;=\; \sum_{t=0}^{W_k-1} q_{s,\,a^{\mathrm{ini}}_k + \delta_d + t,\;d}.
\]

\subsection*{4.3 Decomposição volume $\times$ shape e limite de granularidade}

A separação entre $L_{s,k}$ e $\sigma_{s,t}$ é deliberada: prever o total
de um ciclo é um problema substancialmente mais estável do que prever a
distribuição intra-ciclo, e a equação~(3.2.1) permite que as duas frentes
avancem em paralelo sem sobreposição.

\paragraph{Limitação a declarar.} A base descrita na Seção~4.1 tem
resolução temporal de ciclo (ou, no máximo, de semana, via
\texttt{semana\_pedido}). Ela não tem resolução diária. Como
$\mathrm{Cap}_{c,a}$ é uma capacidade \emph{diária} e $z_{max}$, $z_{min}$
são o pico e o vale \emph{diários} (equações (3.5.2)--(3.5.4)), a
reconstrução de $q_{s,a,d}$ no nível de dia exige uma de três coisas, e o
paper precisa dizer qual foi adotada: (i) uma base diária de outra fonte;
(ii) reformular o modelo em granularidade semanal, com $a$ passando a
indexar semanas e $\mathrm{Cap}$ passando a ser semanal; ou (iii) assumir
um perfil intra-semana por regra de negócio, caso em que $\sigma_{s,t}$ é
em parte premissa, não estimativa.

\subsection*{4.4 Candidatos de modelo}

Todos os candidatos preveem a mesma quantidade, $L_{s,k}$, para os ciclos
do horizonte $K$. A seleção de hiperparâmetros é automática em todos os
casos: com $|S| \approx 800$ séries, ajuste manual por série não é
praticável nem reprodutível.

\begin{center}
\begin{tabular}{@{}p{3.4cm}p{10.2cm}@{}}
\toprule
Candidato & Especificação \\
\midrule
Naive / naive sazonal &
  Baseline obrigatório. Estratégias: último valor; mesmo ciclo do ano
  anterior; média histórica. Serve de referência e de denominador da
  métrica de erro da Seção~4.6. \\
ETS &
  Suavização exponencial em espaço de estados. Seleção automática por AIC
  sobre combinações de erro/tendência/sazonalidade
  $\{A,M\} \times \{N,A,A_d\} \times \{N,A,M\}$, com período sazonal
  $m = \max(\texttt{nm\_ciclo})$. \\
ARIMA / SARIMA &
  Seleção automática de ordem (busca stepwise por AICc) sobre
  $p,q \in \{0,\dots,3\}$, $d \in \{0,1\}$, e ordem sazonal
  $P,D,Q \in \{0,1\}$ com período $m$. \\
SARIMAX &
  Mesma busca, com regressores exógenos derivados da própria base:
  indicadores de mês (\texttt{mes\_pedido}, que captura Dia das Mães,
  Natal e afins) e posição do ciclo no ano (\texttt{nm\_ciclo}). Não exige
  tabela externa de feriados. \\
Prophet &
  \texttt{growth} $\in$ \{linear, logistic\}; \texttt{seasonality\_mode}
  $\in$ \{additive, multiplicative\}; grade em
  \texttt{changepoint\_prior\_scale} $\in \{0{,}01;\,0{,}05;\,0{,}1;\,0{,}5\}$. \\
Gradient boosting &
  LightGBM sobre features de defasagem ($1,2,3,4$ ciclos e ciclo homólogo
  do ano anterior), médias móveis de 3 e 6 ciclos, e os mesmos regressores
  de calendário do SARIMAX. Grade:
  \texttt{n\_estimators} $\in \{100,300,500\}$,
  \texttt{learning\_rate} $\in \{0{,}01;\,0{,}05;\,0{,}1\}$,
  \texttt{num\_leaves} $\in \{15,31,63\}$, com parada antecipada na dobra
  de validação. XGBoost pertence à mesma família e é usado apenas como
  verificação pontual de sensibilidade à biblioteca, não como candidato
  mantido em paralelo. \\
\bottomrule
\end{tabular}
\end{center}

\noindent \textbf{Métodos para demanda intermitente} (Croston, SBA) foram
considerados e descartados: eles se justificam quando a série alvo tem
alta incidência de zeros, o que é típico de contagem de pedidos, não de
volume agregado por ciclo inteiro. Caso a fração de ciclos com
$L_{s,k}=0$ se mostre relevante na base real, o candidato deve ser
reincorporado.

\subsection*{4.5 Nível de ajuste}

A saída exigida pela formulação é por setor individual --- $q_{s,a,d}$ é
indexado por $s$ ---, de modo que não é possível prever apenas em nível
agregado sem desagregar depois.

\paragraph{Agrupamento permitido e proibido.} A combinação Bloco/Subloco
$d$ \textbf{não} pode ser usada como variável de agrupamento ou como
preditor de $L_{s,k}$: $d$ é a variável de decisão do problema
(equação (3.5.1)), e usá-la para explicar o volume introduziria na
previsão justamente aquilo que a otimização pretende alterar, além de
impedir previsões contrafactuais para slots que o setor nunca ocupou. Já
a partição $S_c$ da Seção~3.1 é uma chave de agrupamento legítima, por
ser estável em relação à atribuição.

\paragraph{Estratégia.} Com $|C| = 10$ CDs e $|S| \approx 800$ setores,
$S_c$ oferece grupos de aproximadamente 80 séries cada --- um tamanho
adequado para agrupamento. Adota-se uma estratégia híbrida, comparada
empiricamente pelo critério da Seção~4.6:
\begin{itemize}
  \item modelos estatísticos (ETS, ARIMA/SARIMA, SARIMAX, Prophet)
    ajustados por setor, mas com a busca de hiperparâmetros conduzida por
    grupo $S_c$, de modo que cada setor herde (ou inicialize a busca a
    partir de) a especificação escolhida para seu CD;
  \item gradient boosting ajustado como \textbf{um único modelo global}
    sobre todos os setores simultaneamente, com \texttt{cd\_setor} e o CD
    como variáveis categóricas.
\end{itemize}

\noindent A segunda estratégia segue o resultado da competição M5, em que
modelos treinados conjuntamente sobre muitas séries (\emph{cross-learning})
superaram sistematicamente o ajuste local série a série
\cite{makridakis2022m5}; a literatura de previsão hierárquica registra que
o nível mais desagregado tende a ser o mais ruidoso, o que reforça o
agrupamento \cite{fpp3hierarchical}.

\subsection*{4.6 Protocolo de validação e critério de seleção}

\paragraph{Métrica.} Erro absoluto médio (MAE) e erro absoluto médio
escalonado (MASE) como métricas primárias; RMSE como métrica secundária.
O MASE escalona o erro pelo erro de um benchmark ingênuo, o que o torna
comparável entre setores de portes muito diferentes e o mantém definido
quando $L_{s,k}$ é nulo ou próximo de zero \cite{hyndmanforesight}. O MAPE
é explicitamente rejeitado como critério: com setores de volume baixo, ele
se torna instável ou indefinido. O RMSE é mantido em segundo plano porque
penaliza erros grandes de forma mais acentuada, o que é pertinente dado
que a função objetivo (3.4.1) depende de extremos diários e não de médias.

\paragraph{Validação.} Avaliação por origem deslizante
(\emph{rolling-origin}), nunca por partição aleatória, que violaria a
ordem temporal \cite{fpp3cv}. Para cada origem $k_0$, o modelo é treinado
apenas com ciclos até $k_0$ e avaliado nos $h$ ciclos seguintes, com a
origem deslizando ao longo do histórico e os erros reportados como média
\emph{e} pior dobra. O horizonte $h$ deve ser igual ao número de ciclos
cobertos por $A$ (Bloco~A). A escolha entre janela de treino expansiva e
janela fixa depende do número de anos distintos em \texttt{aa\_ciclo}:
com histórico curto, a janela expansiva é preferível.

\paragraph{Regra de decisão.} Em três níveis:
\begin{enumerate}
  \item menor MASE médio, com peso igual por setor --- e não ponderado por
    volume ---, coerente com o objetivo de nivelamento do problema, que
    trata todos os setores simetricamente;
  \item entre candidatos próximos no critério anterior, menor erro no
    percentil 90 entre setores e dobras: um modelo bom na média mas
    ruim em poucos setores é particularmente nocivo aqui, porque um único
    setor mal previsto pode violar (3.5.4) no CD a que pertence;
  \item persistindo o empate dentro de uma margem pequena, prefere-se o
    modelo mais simples e interpretável.
\end{enumerate}

\subsection*{4.7 Hipóteses e itens em aberto}

\begin{itemize}
  \item \textbf{Invariância de $L_{s,k}$ à atribuição} (Bloco~B.4):
    verificar comparando o volume total de setores comparáveis alocados a
    slots distintos. Se a janela for truncada de forma dependente do dia
    de abertura, a hipótese falha.
  \item \textbf{Esquema de mascaramento} de
    \texttt{total\_volumes\_mascarado}: um fator multiplicativo constante
    é inócuo para a previsão e para o MASE, mas um fator por setor
    distorce tanto o modelo global quanto a agregação por CD da
    restrição (3.5.4).
  \item \textbf{Granularidade efetiva da base}: verificar se
    $(\texttt{cd\_setor}, \texttt{aa\_ciclo}, \texttt{nm\_ciclo})$ é chave
    única. Em caso afirmativo, não há informação intra-ciclo alguma e
    $\sigma_{s,t}$ depende inteiramente de outra fonte; caso contrário, há
    resolução semanal. Verificar também se \texttt{semana\_pedido} é
    semana do mês ou do ano.
  \item \textbf{Dados ausentes} para fechar a formulação: $S_c$,
    $x^{\mathrm{AsIs}}_{s,d}$, $a^{\mathrm{ini}}_k$ e $\mathrm{Cap}_{c,a}$
    (Seção~4.1).
  \item \textbf{Comprimento do histórico}: o número de anos distintos em
    \texttt{aa\_ciclo} determina quantas repetições sazonais existem. Com
    período sazonal da ordem de $\max(\texttt{nm\_ciclo})$, poucos anos de
    histórico tornam a componente sazonal dos modelos SARIMA/ETS pouco
    identificável, o que favorece os modelos globais da Seção~4.5.
\end{itemize}

% =====================================================================
% BLOCO D  -->  SECAO 2 (Related Works): subdividir em 2.1 e 2.2.
% =====================================================================
\section*{Bloco D — §2 Related Works}

A Seção~2 está vazia e hoje cobriria apenas um dos dois problemas do
trabalho. Subdividir em \textbf{2.1 Otimização} (GAP, Simulated Annealing
--- Osman, Metropolis, Kirkpatrick, van Laarhoven e Aarts, já citados na
Seção~4 do paper) e \textbf{2.2 Previsão de demanda}, com o conteúdo
abaixo como ponto de partida:

\begin{quote}
A previsão de demanda em varejo com grande número de séries tem como
referência recente a competição M5, em que os métodos de melhor
desempenho foram, pela primeira vez, integralmente baseados em
aprendizado de máquina --- majoritariamente \emph{gradient boosting} ---,
superando os \emph{benchmarks} estatísticos clássicos e suas combinações;
um dos fatores determinantes foi o \emph{cross-learning}, isto é, o treino
de um único modelo compartilhado entre muitas séries em lugar de um modelo
local por série \cite{makridakis2022m5}. Em estruturas hierárquicas, o
nível mais desagregado é tipicamente o mais ruidoso, o que motiva
estratégias de agregação e reconciliação \cite{fpp3hierarchical}. Quanto à
avaliação, métricas percentuais como o MAPE são inadequadas a séries com
valores nulos ou de baixa magnitude, situação em que se recomendam
métricas escalonadas como o MASE \cite{hyndmanforesight}; e a validação de
modelos de série temporal requer esquemas que preservem a ordem temporal,
como a avaliação por origem deslizante \cite{fpp3cv}.
\end{quote}

\paragraph{Atenção à bibliografia.} No PDF atual, todas as citações
renderizam como \texttt{[?]}: o arquivo \texttt{.bib} não está vinculado
ou está vazio. Isso precisa ser corrigido antes de acrescentar as
referências abaixo, caso contrário apenas se aumenta o número de
\texttt{[?]}. As entradas clássicas já citadas na Seção~4 (Osman;
Metropolis \emph{et al.}; Kirkpatrick \emph{et al.}; Aarts e van
Laarhoven) precisam ser adicionadas com metadados completos --- o Zotero
do projeto resolve isso sem digitação manual.

% =====================================================================
% BLOCO E  -->  SECAO 4.1.5 (que passa a ser 5.1.5): duas frases.
% =====================================================================
\section*{Bloco E — §4.1.5 (que passa a ser §5.1.5)}

Duas inserções curtas na descrição da interface com o simulador:

\begin{enumerate}
  \item No passo 2 (``Agregação da Demanda Diária''), esclarecer a origem
    e o momento do cálculo de $q$:
    \begin{quote}
    Os parâmetros $q_{s,a,d}$ são produzidos pelo módulo de previsão da
    Seção~4, conforme a equação~(3.2.1), e calculados uma única vez antes
    do laço de busca: permanecem fixos ao longo de todas as iterações do
    Simulated Annealing.
    \end{quote}
  \item No parágrafo ``Custo Computacional e Estratégia de
    Implementação'', acrescentar a consequência de estrutura de dados:
    \begin{quote}
    Como $L_{s,k}$ não depende de $d$ (Seção~3.2), a tabela $q_{s,a,d}$
    não precisa ser materializada com $|S| \times |A| \times |D|$
    entradas: basta armazenar $L_{s,k}$ e $\sigma_{s,t}$, com
    $|S| \times \max_k W_k$ entradas, aplicando o deslocamento $\delta_d$
    no momento da consulta. Essa representação é compatível com a
    atualização incremental da carga mencionada acima, já que a
    realocação de um setor corresponde a um deslocamento da sua curva.
    \end{quote}
\end{enumerate}

% =====================================================================
% BLOCO F  -->  arquivo .bib: entradas das referencias novas.
% =====================================================================
\section*{Bloco F — entradas para o \texttt{.bib}}

\begin{verbatim}
@article{makridakis2022m5,
  author  = {Makridakis, Spyros and Spiliotis, Evangelos
             and Assimakopoulos, Vassilios},
  title   = {M5 accuracy competition: Results, findings, and conclusions},
  journal = {International Journal of Forecasting},
  url     = {https://www.sciencedirect.com/science/article/pii/
             S0169207021001874}
  % conferir volume, numero, paginas e ano na fonte antes de submeter
}

@book{fpp3hierarchical,
  author    = {Hyndman, Rob J. and Athanasopoulos, George},
  title     = {Forecasting: Principles and Practice},
  edition   = {3},
  publisher = {OTexts},
  note      = {Cap. 11, Forecasting hierarchical and grouped time series},
  url       = {https://otexts.com/fpp3/hierarchical.html}
}

@book{fpp3cv,
  author    = {Hyndman, Rob J. and Athanasopoulos, George},
  title     = {Forecasting: Principles and Practice},
  edition   = {3},
  publisher = {OTexts},
  note      = {Secao 5.10, Time series cross-validation},
  url       = {https://otexts.com/fpp3/tscv.html}
}

@article{hyndmanforesight,
  author  = {Hyndman, Rob J.},
  title   = {Another look at forecast-accuracy metrics
             for intermittent demand},
  journal = {Foresight: The International Journal of Applied Forecasting},
  url     = {https://robjhyndman.com/papers/foresight.pdf}
  % conferir numero da edicao e ano na fonte antes de submeter
}
\end{verbatim}

\begin{thebibliography}{9}
\bibitem{makridakis2022m5}
Makridakis, S.; Spiliotis, E.; Assimakopoulos, V.
\emph{M5 accuracy competition: Results, findings, and conclusions}.
International Journal of Forecasting.
\url{https://www.sciencedirect.com/science/article/pii/S0169207021001874}

\bibitem{fpp3hierarchical}
Hyndman, R.~J.; Athanasopoulos, G.
\emph{Forecasting: Principles and Practice}, 3.~ed., Cap.~11.
\url{https://otexts.com/fpp3/hierarchical.html}

\bibitem{fpp3cv}
Hyndman, R.~J.; Athanasopoulos, G.
\emph{Forecasting: Principles and Practice}, 3.~ed., Seção~5.10.
\url{https://otexts.com/fpp3/tscv.html}

\bibitem{hyndmanforesight}
Hyndman, R.~J.
\emph{Another look at forecast-accuracy metrics for intermittent demand}.
Foresight: The International Journal of Applied Forecasting.
\url{https://robjhyndman.com/papers/foresight.pdf}
\end{thebibliography}

\end{document}
